\subsection{Primary parity and rational-cycle dependencies}
\label{subsec:primary-parity-dependencies}

The De Bruijn source invokes three earlier results whose exact scopes now
have controlling primary witnesses.  Everett already proves the finite
parity bijection and its nonnegative affine fibres \cite{Everett1977};
Lagarias adds the finite permutation structure and proves the infinite parity
map \cite{Lagarias1985}; Bernstein proves the inverse-series inclusion
\cite{Bernstein1994}, and Lagarias's rational-cycle coordinates
\cite{Lagarias1990}.  This section keeps all three map orientations and both
the raw and reduced denominator layers visible.

Put
\[
 \mathcal R=\Q\cap\Ztwo=\Z_{(2)}.
\]
Lagarias 1985 denotes this ring by $\mathbf Q_2$; it is the set of rationals
with odd reduced denominator, not the full $2$-adic field.  Its elements are
exactly the finite or eventually periodic binary expansions in $\Ztwo$.

\begin{proposition}[Everett's finite bijection; Lagarias's permutation structure]
\label{prop:Lagarias-finite-parity}
For $k\geq1$, define
\[
 \overline Q_k([n]_k)
 =\left[\sum_{i=0}^{k-1}(T^i(n)\bmod2)2^i\right]_k.
\]
Then $\overline Q_k$ is a permutation of $\Z/2^k\Z$.  Its order is a power
of two, in fact it divides $2^k$.  Every permutation cycle of length $r$ at
level $k$ lifts at level $k+1$ either to one cycle of length $2r$ or to two
cycles of length $r$.
\end{proposition}

\begin{proof}
Everett's Theorem~1 proves the bijection by the complete four-case affine
induction in Theorem~\ref{thm:Everett-finite-fibres}.  Lagarias's Theorem B
states the permutation and power-of-two order
\cite[p.~7, Thm.~B]{Lagarias1985}; its proof is explicitly a sketch and omits
the induction details.  Bernstein independently states the lift alternative
\cite[p.~407]{Bernstein1994}.  For the later structure at the scope used here,
Proposition~\ref{prop:BL-Q} gives
\[
 a\equiv b\pmod{2^k}
 \quad\Longleftrightarrow\quad
 Q_3(a)\equiv Q_3(b)\pmod{2^k}.
\]
Thus $\overline Q_k$ is the reduction of the same isometric bijection $Q_3$
and is a permutation.  Each residue has two lifts.  A lifted permutation
cycle either returns each chosen lift after $r$ steps, yielding two
$r$-cycles, or exchanges the two lifts after $r$ steps, yielding one
$2r$-cycle.  Induction from level one proves that every cycle length and the
permutation order divide $2^k$.
\end{proof}

This is the finite labeling called $\Phi_k$ by Laarhoven--de Weger.  Its
cycles are cycles of a permutation of labels.  They are not cycles of the
two-successor modular Collatz relation $\mathcal T_k$, and
Proposition~\ref{prop:Lagarias-finite-parity} alone is not the graph
isomorphism of Theorem~\ref{thm:LDW-finite-debruijn}.

\begin{sourceissue}[Lagarias 1985 finite table and examples]
\label{sourceissue:Lagarias1985-finite-table}
The sentence immediately before Table~2 on printed p.~7 says that one-cycles
are omitted; the table itself is on printed p.~8.  Its $k=6$ row also omits
the nontrivial transposition $(14,46)$.  Direct evaluation gives
\[
 \overline Q_6(14)=46,
 \qquad
 \overline Q_6(46)=14.
\]
The displayed permutation order $4$ remains correct.  The calculation
certificate \path{certificates/lagarias_dependency_checks.py} enumerates
all residues for $k\leq6$ and verifies that this is the only nontrivial cycle
missing from the printed $k=6$ row.

On printed p.~18 the series for $Q_\infty(1)$ has summation index $i$ but
exponent $2n$; the intended exponent is $2i$.  More substantially, the
printed value $Q_\infty(3)=-20/3$ is inconsistent with the paper's own map:
the trajectory $3,5,8,4,2,1,2,\ldots$ has parity word
$11000(10)^\infty$, and hence
\[
 Q_\infty(3)=1+2+\frac{2^5}{1-2^2}=-\frac{23}{3}.
\]
The printed value is not used.
\end{sourceissue}

Lagarias's Eq.~(2.33) defines the infinite map in the forward orientation
\[
 Q_\infty(x)=\sum_{i\geq0}(T^i(x)\bmod2)2^i.
\]

\begin{proposition}[Lagarias infinite parity theorem, proof completed]
\label{prop:Lagarias-infinite-parity}
The map $Q_\infty$ is a continuous, bijective, measure-preserving self-map of
$\Ztwo$.  It is $Q_3$ and Laarhoven--de Weger $\Phi$; it is not the inverse
map $\Phi_3$.
\end{proposition}

\begin{proof}
Theorem L states the result \cite[pp.~17--18]{Lagarias1985}.  Its printed
argument proves prefix preservation, injectivity, surjectivity through finite
inverse residues, and inverse continuity.  Compatibility of those inverse
residues follows because every $\overline Q_k$ is a reduction of the same
prefix map.  Every residue cylinder modulo $2^k$ is therefore permuted onto
one cylinder of the same Haar measure, which supplies the measure-preservation
step not printed at the end of the source proof.
\end{proof}

\begin{proposition}[one proved half of Periodicity]
\label{prop:Bernstein-periodicity-half}
Let $\Phi_3=Q_3^{-1}$.  Then
\begin{equation}
\label{eq:Bernstein-rational-inclusion}
 \Phi_3(\mathcal R)\subseteq\mathcal R,
\end{equation}
or equivalently
\begin{equation}
\label{eq:Bernstein-rational-inclusion-forward}
 \mathcal R\subseteq Q_3(\mathcal R).
\end{equation}
The Periodicity Conjecture is the equality
\begin{equation}
\label{eq:Lagarias-periodicity-conjecture}
 Q_3(\mathcal R)=\mathcal R.
\end{equation}
\end{proposition}

\begin{proof}
Bernstein writes a rational parity address as a finite or eventually
periodic sequence of one-positions and substitutes it into his explicit
series for $\Phi_3$.  The resulting geometric relation proves that the image
is rational \cite[p.~407, Cor.~1]{Bernstein1994}.  Applying the bijection
$Q_3$ gives the equivalent forward inclusion.  Lagarias prints
\eqref{eq:Lagarias-periodicity-conjecture} as an unproved conjecture on
printed p.~18 \cite{Lagarias1985}; no reverse inclusion is added.
\end{proof}

The positive-integer ``third form'' in Lagarias 1985 is likewise a conjecture:
\begin{equation}
\label{eq:Lagarias-third-form}
 Q_3(\Npos)\subseteq\frac13\Z\setminus\Z.
\end{equation}
Bernstein instead proves that the inverse-oriented assertion
$\Z^+\subseteq\Phi_3((1/3)\Z)$ is equivalent to the ordinary conjecture
\cite[pp.~406--407, Thms.~2--3]{Bernstein1994}.  An equivalence to an endpoint
conjecture is not a proof of either formulation.

Lagarias's last exploratory value on printed p.~18 is left as
$Q_3(-82/7)=?/15$.  Exact iteration gives a prefix of length $13$ followed by
the parity block $1100$, hence the independent finite calculation
\begin{equation}
\label{calc:Lagarias1985-minus82}
 Q_3(-82/7)=\frac{18098}{5}=\frac{54294}{15}.
\end{equation}
The same certificate reproduces this value.  It is not attributed to the
source and does not prove Periodicity.

\begin{theorem}[Lagarias rational fixed-return coordinates]
\label{thm:Lagarias-rational-return}
Let $n\geq1$ and $v=(v_0,\ldots,v_{n-1})\in\{0,1\}^n$, and put
\[
 m(v)=\sum_{i=0}^{n-1}v_i,
 \qquad
 D(v)=2^n-3^{m(v)}\neq0,
\]
\[
 N(v)=\sum_{j=0}^{n-1}
 v_j2^j3^{v_{j+1}+\cdots+v_{n-1}}.
\]
There is exactly one $x(v)\in\mathcal R$ with
$T^n(x(v))=x(v)$ and initial parity word $v$, namely
\begin{equation}
\label{eq:Lagarias-rational-return}
 x(v)=\frac{N(v)}{D(v)}.
\end{equation}
Its least $T$-period is $n$ exactly when $v$ is not a proper repeated block.
\end{theorem}

\begin{proof}
The branch composition determined by $v$ is
\[
 U_{v_{n-1}}\circ\cdots\circ U_{v_0}(x)
 =\frac{3^{m(v)}x+N(v)}{2^n},
\]
and $D(v)\neq0$ because $2^n$ is even while $3^{m(v)}$ is odd.  Thus its
unique formal fixed point is
\eqref{eq:Lagarias-rational-return}.  Lagarias proves that this formal point
has the prescribed actual branches: for each $y\in\mathcal R$, exactly one of
the two forward branch expressions $U_0(y)=y/2$ and
$U_1(y)=(3y+1)/2$ lies in $\mathcal R$, with the choice determined by the
parity of $y$; if $y\notin\mathcal R$, both expressions remain outside
$\mathcal R$ \cite[pp.~36--37, Thm.~2.1]{Lagarias1990}.  The denominator
$D(v)$ is odd, so the formal fixed point lies in $\mathcal R$.  If its
prescribed composition ever selected the other forward expression, it would
leave $\mathcal R$ and could never return, contradicting that the composition
ends again at $x(v)\in\mathcal R$.  Thus every prescribed branch is actual.
For every $d\geq1$, the conjugacy in
Proposition~\ref{prop:BL-Q} gives
$Q_3(T^d(x(v)))=S^d(Q_3(x(v)))$.  Hence a $T$-return after $d$ steps implies
a shift return after $d$ steps.  Conversely, a shift return and the
injectivity of $Q_3$ from Proposition~\ref{prop:Lagarias-infinite-parity}
imply $T^d(x(v))=x(v)$.  The two sets of positive return times are equal, so
their least elements agree; this least period is $n$ exactly when $v$ is not
a proper repeated block.
\end{proof}

Write the positive reduced common denominator as
\begin{equation}
\label{eq:Lagarias-reduced-denominator}
 \delta(v)=\frac{|D(v)|}{\gcd(|N(v)|,|D(v)|)}.
\end{equation}
It is always odd and prime to $3$.  The proof cannot omit the all-even word.
If $m(v)\geq1$, then $3\nmid D(v)$.  If $m(v)=0$, the return equation is
$x=x/2^n$, so $x=0$ and $\delta(v)=1$.  Thus
\begin{equation}
\label{eq:Lagarias-denominator-mod-six}
 \delta(v)\equiv\pm1\pmod6.
\end{equation}

\begin{proposition}[raw scaling fibres and the primitive layer]
\label{prop:Lagarias-primitive-scaling}
For odd $b$, let
\[
 T^{(3,b)}(n)=
 \begin{cases}
 (3n+b)/2,&n\text{ odd},\\
 n/2,&n\text{ even}.
 \end{cases}
\]
Multiplication by a common positive odd denominator gives
\[
 bT(x_j)=T^{(3,b)}(bx_j).
\]
Consequently rational $T$-cycles correspond to common-odd-scaling classes
of integer $T^{(3,b)}$-cycles.  Without identifying the raw representatives,
the unique member with $b$ equal to the positive reduced common denominator
is characterized by
\[
 \gcd(b,n_0,\ldots,n_{L-1})=1.
\]
It is exactly a primitive integer $T^{(3,b)}$-cycle, and division by $b$ is
the two-sided inverse.  Its parameter satisfies $b\equiv\pm1\pmod6$.
\end{proposition}

\begin{proof}
The branch identity is the calculation in
Proposition~\ref{prop:LDW-rational-cycle-scaling}.  Lagarias explicitly
defines the primitive condition and the inverse correspondence on printed
pp.~39--40 \cite{Lagarias1990}.  For the raw strata, if $d\mid b$ and
$n=dm$, $b=dc$, then
\[
 T^{(3,b)}(dm)=dT^{(3,c)}(m).
\]
This proves every scaling fibre without suppressing it.  Choosing the reduced
common denominator gives the displayed gcd condition and a literal inverse.
Equation~\eqref{eq:Lagarias-denominator-mod-six} supplies the parameter
restriction.
\end{proof}

\begin{sourceissue}[Lagarias 1990 local coordinate defects]
\label{sourceissue:Lagarias1990-local-defects}
Six printed statements require exact local repair.

First, the itinerary sentence at the top of printed p.~34 defines
$v=(v_0,\ldots,v_{n-1})$ but quantifies
$T^i(x(v))\equiv v_i\pmod2$ for $1\leq i\leq n$, which skips $v_0$ and invokes
the undefined $v_n$.  The correct range is $0\leq i\leq n-1$.  Second,
printed p.~39 says that $T$ never changes the denominator of an
arbitrary element of $\mathcal R$, but $T(1/3)=1$.  The correct preserved
scope is a reduced odd denominator prime to $3$; every periodic rational has
that property by the case split above.  Third, the printed formula for
$\delta(v)$ needs $|D(v)|$ for a positive denominator unless a signed-gcd
convention is supplied.  Fourth, printed p.~40 writes
$T^{(3,b)}(n)=T^{(3,-b)}(-n)$, whereas direct branch evaluation gives
\[
 T^{(3,b)}(n)=-T^{(3,-b)}(-n).
\]
Fifth, printed p.~34 ends one correspondence sentence with
$b\equiv1\pmod6$; the preceding sentence, printed p.~40, and the examples
require $b\equiv\pm1\pmod6$.  Sixth, the invariant gcd stratum on p.~40 uses
the standing $\gcd(b,3)=1$ hypothesis and is not a statement for arbitrary
odd $b$.
\end{sourceissue}

\begin{nonclaim}
The Periodicity Conjecture does not follow from the one-sided Bernstein
inclusion.  Lagarias's primitive-cycle existence and finiteness statements
are conjectures.  A periodic parity address gives a rational periodic state,
but this does not classify the positive integers inside $\Ztwo$, exclude
nontrivial integer cycles, or prove convergence to $1$.
\end{nonclaim}
